\chapter{Perspectives}

La prochaine phase du projet doit maintenant passer d'une reproduction
fonctionnelle credible a une etude critique plus ambitieuse. Trois axes se
degagent.

\section{Fidelite supplementaire}

Le premier axe consiste a rapprocher encore l'implementation du protocole des
auteurs. Les priorites identifiees sont l'amelioration du couple coherence /
naming, une meilleure maitrise de la recherche de \texttt{K}, et une reduction
de la sur-fragmentation des clusters finaux.

\section{Evaluations complementaires}

Le deuxieme axe consiste a produire des evaluations qui depassent la seule
reproduction du papier :
\begin{itemize}[leftmargin=*]
  \item ablations plus systematiques ;
  \item nouveaux jeux de donnees ;
  \item variations d'hyperparametres ;
  \item comparaisons avec des approches plus simples ou plus classiques.
\end{itemize}

\section{Ameliorations proposees}

Le troisieme axe vise a transformer le projet en contribution personnelle.
L'objectif n'est plus seulement de refaire le papier, mais de proposer des
ameliorations algorithmiques ou computationnelles evaluables, puis de mesurer
leur impact de maniere rigoureuse.
